% Part: normal-modal-logic
% Chapter: sequent-calculus
% Section: rules-for-K

\documentclass[../../../include/open-logic-section]{subfiles}

\begin{document}

\olfileid[pt]{nml}{seq}{rul}

\olsection{Regras para \Log{K}}

As regras para os conectivos proposicionais usuais são as mesmas do
cálculo de sequentes usual~$\Log{LK}$. Os axiomas também são os mesmos:
todo sequente da forma $!A \Sequent !A$ conta como axioma.

Para a linguagem modal, com\iftag{prvBox}{\iftag{prvDiamond}{~$\Box$
e}{~$\Box$}}{}\iftag{prvDiamond}{~$\Diamond$}{}, temos a
\iftag{notprvBox,notprvDiamond}{seguinte regra adicional}{seguintes regras adicionais}:
\[
  \iftag{prvBox}{\iftag{prvDiamond}{% <> and [] primitive
    \Axiom$\Gamma \fCenter \Delta, !A$
    \RightLabel{$\Box$}
    \UnaryInf$\Box\Gamma \fCenter \Diamond\Delta, \Box!A$
    \DisplayProof  
    \qquad
    \Axiom$!A, \Gamma \fCenter \Delta$
    \RightLabel{$\Diamond$}
    \UnaryInf$\Diamond!A, \Box\Gamma \fCenter \Diamond\Delta$
    \DisplayProof
}{% only Box primitive
\Axiom$\Gamma \fCenter !A$
\RightLabel{$\Box$}
\UnaryInf$\Box\Gamma \fCenter \Box!A$
\DisplayProof
}}{% only <> primitive
  \Axiom$!A \fCenter \Delta$
  \RightLabel{$\Diamond$}
  \UnaryInf$\Diamond!A \fCenter \Diamond\Delta$
  \DisplayProof
}
\]
Aqui, \iftag{prvBox}{$\Box\Gamma$ significa a sequência de !!{formula}s
obtida de~$\Gamma$ antepondo $\Box$ a cada
!!{formula} de~$\Gamma$\iftag{prvDiamond}{ e
}{}}{}\iftag{prvDiamond}{$\Diamond\Delta$ é a sequência de
!!{formula}s obtida de~$\Delta$ antepondo $\Diamond$ a cada
!!{formula} de~$\Delta$}{}. \iftag{notprvDiamond}{O lado direito da
premissa da regra $\Box$ deve conter no máximo uma
!!{formula}~$!A$. }{}%
\iftag{notprvBox}{O lado esquerdo da premissa da regra $\Diamond$
deve conter no máximo uma !!{formula}~$!A$. }{}%
\iftag{prvBox}{$\Gamma$\iftag{prvDiamond}{~e~}{}}{}\iftag{prvDiamond}{$\Delta$}{}
podem ser vazios; nesse caso, a parte correspondente
\iftag{prvBox}{$\Box\Gamma$\iftag{prvDiamond}{~e~}{}}{}\iftag{prvDiamond}{$\Diamond\Delta$}{}
do sequente da conclusão também é vazia.

A restrição de acrescentar
\iftag{prvBox}{$\Box$ à direita\iftag{prvDiamond}{ e
}{}}{}\iftag{prvDiamond}{$\Diamond$ à esquerda}{}
a uma única !!{formula}~$!A$ é necessária. Se fosse permitido
\iftag{prvBox}{acrescentar $\Box$ a qualquer número de !!{formula}s à
direita\iftag{prvDiamond}{ ou }{}}{}\iftag{prvDiamond}{acrescentar $\Diamond$
a qualquer número de !!{formula}s à esquerda}{}, seria possível
!!{derive}:
  \[
    \iftag{prvBox}{
      \Axiom$!A \fCenter !A$
      \RightLabel{\RightR{\lnot}}
      \UnaryInf$\fCenter !A, \lnot !A$
      \RightLabel{$\Box*$}
      \UnaryInf$\fCenter \Box!A, \Box\lnot !A$
      \doubleLine
      \RightLabel{\RightR{\lor}}
      \UnaryInf$\fCenter \Box!A \lor \Box \lnot !A$
      \DisplayProof}{}
    \iftag{notprvBox,notprvDiamond}{}{\qquad}
    \iftag{prvDiamond}{
      \Axiom$!A \fCenter !A$
      \RightLabel{\LeftR{\lnot}}
      \UnaryInf$\lnot !A, !A \fCenter$
      \RightLabel{$\Diamond*$}
      \UnaryInf$\Diamond\lnot !A,\Diamond!A \fCenter $
      \RightLabel{\RightR{\lnot}}
      \UnaryInf$\Diamond!A \fCenter \lnot\Diamond\lnot !A$
      \RightLabel{\RightR{\lif}}
      \UnaryInf$ \fCenter \Diamond!A \lif \lnot\Diamond\lnot !A$
      \DisplayProof}{}
  \]
Mas \iftag{prvBox}{$\Box!A \lor \Box \lnot !A$\iftag{prvDiamond}{ e
$\Diamond!A \lif \lnot\Diamond\lnot !A$ não são válidas}{ não é válida}}{$\Diamond!A \lif
\lnot\Diamond\lnot !A$ não é válida} em~$\Log{K}$.

Se fossem permitidas fórmulas laterais além de~$!A$ na premissa e se
permitisse \iftag{prvBox}{que a regra $\Box$ acrescentasse $\Box$ apenas a~$!A$ à
direita\iftag{prvDiamond}{, ou se permitisse }{}}{}\iftag{prvDiamond}{que a
regra $\Diamond$ acrescentasse $\Diamond$ apenas a~$!A$ à esquerda}{} (sem
fazer nada com as fórmulas laterais), seria possível !!{derive}:
\[
  \iftag{prvBox}{
    \Axiom$!A \fCenter !A$
    \RightLabel{\RightR{\lnot}}
    \UnaryInf$\fCenter !A, \lnot !A$
    \RightLabel{\RightR{\Exchange}}
    \UnaryInf$\fCenter \lnot !A, !A$
    \RightLabel{$\Box*$}
    \UnaryInf$\fCenter \lnot!A, \Box !A$
    \doubleLine
    \RightLabel{\RightR{\lor}}
    \UnaryInf$\fCenter \lnot!A \lor \Box !A$
    \DisplayProof}{}
  \iftag{notprvBox,notprvDiamond}{}{\qquad}
  \iftag{prvDiamond}{
    \Axiom$!A \fCenter !A$
    \RightLabel{\LeftR{\lnot}}
    \UnaryInf$\lnot !A, !A \fCenter$
    \RightLabel{$\Diamond*$}
    \UnaryInf$\Diamond\lnot !A, !A \fCenter $
    \RightLabel{\RightR{\lnot}}
    \UnaryInf$!A \fCenter \lnot\Diamond\lnot !A$
    \RightLabel{\RightR{\lif}}
    \UnaryInf$ \fCenter !A \lif \lnot\Diamond\lnot !A$
    \DisplayProof}{}
\]
Mas \iftag{prvBox}{$\lnot!A \lor \Box !A$ (que é equivalente a $!A
\lif \Box!A$)\iftag{prvDiamond}{ e $!A \lif \lnot\Diamond\lnot !A$
não são válidas}{ não é válida}}{$!A \lif \lnot\Diamond\lnot !A$ não é válida} em~$\Log{K}$.

\end{document}
